%% P16 --- Jacobians, the chain rule and automatic differentiation
%%
%% Every number in this program that the reader cannot do in their head comes
%% from code/p16_autodiff.py and is pulled in with \val{}. The console listing
%% is a committed transcript written by that script.
%%
%% THE THESIS: which mode is cheaper is a question about SHAPE and nothing
%% else. The chain rule multiplies Jacobians; multiplication is associative
%% but its COST is not, which Program P06 measured; so the two ways of
%% bracketing the same product are forward and reverse mode, and a scalar loss
%% with many parameters makes one cheaper by a factor of the parameter count.
%%
%% WHAT THIS PROGRAM IS OWED, read out of the files rather than remembered:
%%   F12  gives the chain rule and the product over layers, and says of sigma'
%%        being expressible in sigma that it "is the seed of what Program P16
%%        calls reverse mode". It also names the two things this program owes
%%        arithmetic for: the activations must be stored, and checkpointing
%%        trades recomputation for memory.
%%   P06  MEASURED THE BRACKETING already -- one triple product, two
%%        bracketings, identical answers, a factor of 500 in cost. That IS
%%        this program's result one program early, so the ratio is GATED
%%        against P06's committed values rather than re-derived.
%%   P13  gives the DAG and the theorem that every topological order computes
%%        the same values, which is what lets reverse mode walk backwards.
%%   P15  gives the gradient, which is the one-row Jacobian.
%%   P03  already NAMES activations and checkpointing and says the trade is
%%        operations for bytes. It does not do the arithmetic. This program
%%        owes exactly that and must not restate the concept.
%%
%% WHAT IT LEAVES ALONE, checked against tools/programs.json:
%%   the Hessian, curvature and the step-size bound               -> P17
%%   matrix calculus, the softmax-cross-entropy gradient          -> P18
%%   the optimisers                                               -> P20
%%   the residual stream                                          -> P32
%%
%% NO LIBRARY IS NAMED, on F04's rule that a framework's internals are a fact
%% about a version. What is here is a reverse-mode differentiator in about
%% forty lines, which exhibits every claim including all three silent
%% failures -- much more convincing demonstrated than asserted.
%%
%% Twin: programs/pl/P16-autodiff.tex. Frame for frame, macro for macro,
%% number for number; tools/parity.py compares them.

\program[Jacobians and automatic differentiation]{Jacobians, the chain rule
  and automatic differentiation}
\label{prog:P16}
\index{autodiff}
\index{Jacobian}

\begin{outcomes}
\outcome{1--8}{say what shape the derivative of a vector-valued function is,
  and why the gradient is a special case of it}
\outcome{9--15}{say what the chain rule does to those matrices, and why
  nobody ever forms one}
\outcome{16--22}{say which of the two modes is cheaper for a given problem,
  from its shape alone}
\outcome{23--29}{say what reverse mode has to keep, what that costs, and what
  recomputing some of it buys}
\outcome{30--38}{name three situations in which the number you get back is
  not the derivative you meant, and nothing says so}
\end{outcomes}

\begin{quiz}
\item A function takes $m$ inputs and returns $n$ outputs. What shape is its
  derivative? \teachesat{1--2}
  \answerto{An $n \times m$ matrix \dash{} one row per output, one column per
  input. It is called the Jacobian, and Program~\ref{prog:P15}'s gradient is
  the case $n = 1$.}
\item Why does a framework not compute the Jacobian of a network?
  \teachesat{13--14}
  \answerto{Because it would have one row per output and one column per
  parameter, which for a real model is a matrix nobody can store. What is
  computed is a \emph{product} with it, which is the same size as the
  parameters and no larger.}
\item Training uses reverse mode rather than forward mode. What fact about
  the problem decides that? \teachesat{17--18}
  \answerto{Its shape, and nothing else: many inputs, one output. Forward mode
  costs one pass per input and reverse one per output, so with a scalar loss
  reverse wins by a factor of the parameter count.}
\item Why does reverse mode need more memory than a forward pass alone?
  \teachesat{23--24}
  \answerto{Because working backwards needs each layer's own input, and those
  were computed on the way forward. They have to be kept, so memory grows with
  depth and with batch size.}
\item Your gradient check passes to four decimal places and not to fifteen.
  Is the implementation wrong? \teachesat{28--29}
  \answerto{Almost certainly not. A finite difference is a testing instrument
  with an error floor of its own \dash{} Program~\ref{prog:F11} measured that
  floor \dash{} so agreement to a handful of digits is what a correct
  implementation looks like.}
\end{quiz}

% ============================================================
\section{What shape a derivative has}
\label{sec:P16-jacobian}
\index{Jacobian}

\begin{fr}
Program~\ref{prog:P15} differentiated a function of several inputs that
returned one number, and collected the partial derivatives into a gradient.

A layer of a network does not return one number. It takes a vector in and
returns a vector out. So the question this program starts from is the obvious
next one: if $f$ takes $m$ inputs and returns $n$ outputs, how many partial
derivatives are there, and how would you arrange them?
\dotline
\nextframe
\end{fr}

\begin{fr}
\begin{ansblock}
There are $n \times m$ of them \dash{} one for every output-input pair
\dash{} and they arrange as a matrix with one row per output and one column
per input.
\end{ansblock}

That matrix is the \textbf{Jacobian}, written $J$ or
$\frac{\partial f}{\partial x}$:
\[ J_{ij} = \frac{\partial f_i}{\partial x_j} \]
Row $i$ is the gradient of the $i$th output. Column $j$ says what happens to
everything when input $j$ is nudged.

\mermaidfig{p16-jacobian-shape}{Rows are outputs, columns are inputs.}{the
shape of a derivative}
\end{fr}

\begin{fr}
Check that against what you already have. What is the Jacobian of a function
with exactly one output?
\nextframe
\end{fr}

\begin{fr}
\begin{ansblock}
A single row \dash{} which is Program~\ref{prog:P15}'s gradient.
\end{ansblock}

So the gradient is not a different object; it is the Jacobian of the case that
matters most, and every sentence P15 proved about it is a sentence about a
one-row matrix.

\begin{note}
There is a transpose convention here that costs people an afternoon. The
gradient is usually written as a \emph{column} vector, and the Jacobian's
single row is a \emph{row} vector, so the two differ by a transpose.

Nothing mathematical turns on it, and this book will not fuss over it, but it
is worth knowing that is where the disagreement between two textbooks usually
lives. Program~\ref{prog:P18} has to be careful about it because it
differentiates matrix expressions; here it changes nothing.
\end{note}
\end{fr}

\begin{fr}
One more, because it is the case a whole layer is. What is the Jacobian of a
linear map $f(x) = Ax$?
\dotline
\nextframe
\end{fr}

\begin{fr}
\ans{$A$ itself}
Because $f_i = \sum_j A_{ij} x_j$, so $\partial f_i / \partial x_j$ is
$A_{ij}$, which is the definition of $A$.

\textbf{A linear map is its own derivative}, which is exactly what
Program~\ref{prog:F11} said about a straight line in one dimension and is the
same statement one dimension count up. It is also why the Jacobian of a
network's linear layer costs nothing to think about: it is the weight matrix,
already in memory.

\begin{aibox}
The interesting Jacobians in a network are the ones that are not weight
matrices. An elementwise activation has a \emph{diagonal} Jacobian, because
output $i$ depends on input $i$ and on nothing else \dash{} so its $n^{2}$
entries are $n$ numbers and $n^{2} - n$ zeros.

That is the first hint of this program's whole argument: the useful thing
about a Jacobian is almost never its entries. It is what multiplying by it
does, and how cheaply.
\end{aibox}
\end{fr}

\begin{fr}
Before going on, one sanity check on size. A network with
$\val{p16.params}$ parameters and a scalar loss. What are the dimensions of
the Jacobian of the loss with respect to the parameters, and how many numbers
is that?
\nextframe
\end{fr}

\begin{fr}
\begin{ansblock}
One row and $\val{p16.params}$ columns \dash{} so it is the gradient, and it
is the same size as the model.
\end{ansblock}

Which is large but is exactly what you expected, because it is one number per
parameter and you already knew you were going to get that.

Hold that thought, because the next section is about a Jacobian that is
\emph{not} the same size as the model, and about what is done instead of
forming it.
\end{fr}

% ============================================================
\section{The chain rule multiplies them, and nobody forms one}
\label{sec:P16-chain}
\index{autodiff!chain rule}

\begin{fr}
Program~\ref{prog:F12} gave the chain rule for one variable: differentiate the
outside, multiply by the derivative of the inside.

For vector-valued functions the same sentence holds with one word changed.
If $y = g(f(x))$, then
\[ J_{g \circ f}(x) = J_g\bigl(f(x)\bigr)\; J_f(x) \]
\dash{} the Jacobians \emph{multiply}, in that order.

Why that order rather than the other? The answer is in
Program~\ref{prog:P06}.
\dotline
\nextframe
\end{fr}

\begin{fr}
\begin{ansblock}
Because Program~\ref{prog:P06} showed matrix multiplication is composition,
and the shapes only fit one way round.
\end{ansblock}

$J_f$ maps the input space to the middle space; $J_g$ maps the middle space to
the output space; so $J_g J_f$ is a map from input to output, and $J_f J_g$ is
usually not even defined. The chain rule is composition of linear maps and
nothing more.

\textbf{So a network of $L$ layers has a Jacobian that is a product of $L$
matrices}, one per layer, and this is Program~\ref{prog:F12}'s product over
layers with the word \enquote{number} replaced by \enquote{matrix} throughout.
\end{fr}

\begin{fr}
Now the practical question, and it is the one that turns this from notation
into engineering.

Take a layer with $\val{p16.width}$ inputs and $\val{p16.width}$ outputs, in
a stack $\val{p16.depth}$ deep. How big is each Jacobian, and what would it
cost to form the whole product?
\dotline
\nextframe
\end{fr}

\begin{fr}
\begin{ansblock}
Each is $\val{p16.width} \times \val{p16.width}$, and forming the product means
multiplying $\val{p16.depth}$ of them together.
\end{ansblock}

At those sizes that is nothing. Now put a real layer in: a width of a few
thousand rather than $\val{p16.width}$, and a depth in the tens.

\textbf{Each Jacobian is then a matrix of millions of entries, and there is one
per layer.} Nobody stores that, and nobody needs to \dash{} because the answer
you actually want is not the product. It is the product \emph{applied to
something}.
\end{fr}

\begin{fr}
That is the sentence the rest of the program turns on, so it is worth being
precise about it.

You never want $J$. You want either $Jv$ for some direction $v$, or $u^{\mathsf{T}}J$
for some weighting $u$. Both of those are vectors, not matrices.

Given a chain $J_3 J_2 J_1$ and a vector $v$, write down two different ways to
bracket the computation of $J_3 J_2 J_1 v$.
\dotline
\nextframe
\end{fr}

\begin{fr}
\begin{ansblock}
$J_3(J_2(J_1 v))$, working from the right; or $(J_3 J_2) J_1$ and then apply
it to $v$, working from the left.
\end{ansblock}

Both give the same answer, because matrix multiplication is associative. They
do not cost the same, and Program~\ref{prog:P06} measured exactly that: one
bracketing of a triple product cost $\val{p06.cost.left}$ multiplications and
the other $\val{p06.cost.right}$, a factor of $\val{p06.cost.ratio}$, for an
identical result.

\textbf{That measurement is this program's whole subject, taken one program
before it had a name.} The right-to-left bracketing is forward mode. The
left-to-right bracketing is reverse mode. The script gates the ratio here
against P06's committed numbers, so the two programs cannot come apart.

\mermaidfig{p16-two-brackets}{The same product, from either end.}{two
brackets}
\end{fr}

\begin{fr}
\begin{note}
Notice what the bracketing does to the shapes, because that is the whole
mechanism. Working from the right, every intermediate is a \emph{vector}: a
matrix times a vector is a vector, so nothing bigger than a layer's worth of
numbers ever exists. Working from the left with a row vector, the same is true
from the other end.

Working from the middle outwards would form a matrix, and that is the thing
neither mode ever does. \textbf{Nobody forms a Jacobian} is not a
simplification; it is a statement about which brackets are used.
\end{note}
\end{fr}

% ============================================================
\section{Which end, and why the shape decides}
\label{sec:P16-modes}
\index{autodiff!forward and reverse mode}

\begin{fr}
Name the two directions properly, because the names are used constantly and
they say what they do.

\textbf{Forward mode} starts at the input end. It pushes one direction $v$
through the whole chain, and what comes out is $Jv$ \dash{} how every output
moves when you nudge the inputs along $v$.

\textbf{Reverse mode} starts at the output end. It pulls one weighting $u$
back through the whole chain, and what comes out is $u^{\mathsf{T}}J$ \dash{}
how one weighted combination of outputs depends on every input.

One pass of forward mode gives you one \emph{column's worth}. One pass of
reverse gives you one \emph{row's worth}. If you wanted the whole Jacobian, how
many passes would each need?
\dotline
\nextframe
\end{fr}

\begin{fr}
\begin{ansblock}
Forward: one per input. Reverse: one per output.
\end{ansblock}

And there is the entire decision procedure. It has nothing to do with the
network, the data, the framework or the hardware. It is two integers.

\textbf{Many inputs and few outputs: reverse. Few inputs and many outputs:
forward.}

Now apply it. Training a model computes the derivative of one scalar loss with
respect to every parameter. Which mode, and by what factor?
\nextframe
\end{fr}

\begin{fr}
\begin{ansblock}
Reverse, by a factor of the number of parameters.
\end{ansblock}

One output means one reverse pass. $\val{p16.params}$ inputs would mean
$\val{p16.params}$ forward passes. The factor is not large, it is absurd: it
is the size of the model.

\textbf{That is why every training loop in existence uses reverse mode, and it
is a consequence of the loss being a single number.} Nothing else about deep
learning is involved in the argument.

\begin{aibox}
It is worth noticing what would change the answer. A function with one input
and many outputs \dash{} a sensitivity analysis over a single parameter, say
\dash{} is forward mode's case, and using reverse there would be the same
absurdity in the other direction.

Frameworks provide both for that reason. The reason you have only ever used
one of them is that training has exactly one shape.
\end{aibox}
\end{fr}

\begin{fr}
The measurement, on a chain small enough to count exactly. A stack
$\val{p16.depth}$ deep and $\val{p16.width}$ wide, ending in one output, with
every Jacobian a rational matrix so nothing rounds.

Both routines count their own multiplications as they go. Predict the ratio
before reading on.
\dotline
\nextframe
\end{fr}

\begin{fr}
\ans{$\val{p16.mode.ratio}$, which is the number of inputs}
$\val{p16.fwd.muls}$ multiplications forward against
$\val{p16.rev.muls}$ reverse, for a gradient that is \emph{identical} \dash{}
asserted exactly over the rationals, because a difference between the two
would have made the comparison about something else.

And the ratio is not a fact about this network. The script sweeps widths of
$3$, $7$, $\val{p16.width}$ and $50$ and asserts the ratio equals the width
every time, which is the invariant; any single figure would have been a fact
about one stack.

\begin{trapbox}
\textbf{\enquote{Reverse mode is faster.}} It is faster \emph{for this shape},
and the shape is the whole reason. Said without the qualifier it becomes
folklore, and folklore is what makes people reach for it on a problem with one
input and ten thousand outputs, where it is exactly as wasteful as forward
mode is here.

There is a second half people drop, and section 4 is about it: reverse mode
buys its speed with memory, and forward mode does not need any.
\end{trapbox}
\end{fr}

\begin{fr}
Before that, the connection to Program~\ref{prog:P13}, because it is what
makes the whole thing well defined.

A network's computation is a directed graph with no cycles: every value is
produced before it is used. Reverse mode walks that graph \emph{backwards},
accumulating a derivative at each node from the ones that consume it.

What does Program~\ref{prog:P13} guarantee that makes this legitimate?
\nextframe
\end{fr}

\begin{fr}
\begin{ansblock}
That the graph has a topological order at all, and that every one of them
computes the same values.
\end{ansblock}

So the backward walk may visit the nodes in any order that respects the
dependencies, and it will get the same answer. That is why an implementation
is free to reorder, to parallelise and to fuse without anyone checking the
arithmetic again.

\begin{note}
It also says what breaks it: a cycle. Program~\ref{prog:P13} showed that a
cycle does not make evaluation slow, it makes it \emph{undefined} \dash{}
there is no order in which each node follows what it needs.

A recurrent network looks like a cycle and is not: it is unrolled into a DAG
before anything is differentiated, which is why its memory grows with the
number of steps unrolled. That is P13's own observation and it is the reason
the sentence \enquote{backpropagation through time} contains the word
\emph{through}.
\end{note}
\end{fr}

% ============================================================
\section{What it has to keep}
\label{sec:P16-memory}
\index{autodiff!activations}

\begin{fr}
Reverse mode's cost is not arithmetic, and this is the half that surprises
people.

Look at one layer's Jacobian in a chain. To multiply by it on the way back,
you need its entries. Its entries depend on that layer's \emph{input}, which
was computed on the way forward.

What follows about memory?
\dotline
\nextframe
\end{fr}

\begin{fr}
\begin{ansblock}
Every layer's input has to be kept from the forward pass until the backward
pass reaches it. Memory grows with depth, and again with batch size.
\end{ansblock}

Those kept values are the \textbf{activations}, and Program~\ref{prog:P03}
already counted them in its memory bill: they do not scale with the parameter
count, they scale with depth times batch times width, so the same model has a
different activation bill on every workload.

\textbf{Forward mode needs none of this.} It carries one direction alongside
the computation and never looks back, so its memory is a constant. That is the
second half of the trade the previous section's trapbox named: reverse mode
buys a factor of the parameter count in time and pays for it in memory.

\mermaidfig{p16-keep-or-recompute}{What the backward pass needs, and the way
out.}{keep or recompute}
\end{fr}

\begin{fr}
Program~\ref{prog:P03} also named the way out \dash{} gradient checkpointing,
which keeps only some of the activations and recomputes the rest \dash{} and
deliberately did not do the arithmetic. Here it is.

Keep every $k$th activation in a stack $L$ deep. Between two kept ones you
must recompute a segment before you can walk back through it. Write down the
peak number of activations held at once.
\dotline
\nextframe
\end{fr}

\begin{fr}
\begin{ansblock}
$\frac{L}{k} + k$: the $\frac{L}{k}$ checkpoints, plus the $k$ activations of
the one segment being recomputed.
\end{ansblock}

That is a function with a minimum, and it is worth finding rather than
guessing. Differentiate it: $-\frac{L}{k^{2}} + 1$, which is zero at
$k = \sqrt{L}$.

\textbf{So the best checkpoint spacing is the square root of the depth, and the
peak memory is $2\sqrt{L}$.} At $L = \val{p16.ckpt.layers}$ that is a
checkpoint every $\val{p16.ckpt.bestk}$ layers and
$\val{p16.ckpt.peak}$ activations held at once, against
$\val{p16.ckpt.layers}$ \dash{} a factor of $\val{p16.ckpt.saving}$.

And the cost is exactly $\val{p16.ckpt.passes}$ forward passes rather than one,
because each segment is recomputed exactly once.

\begin{note}
The $\sqrt{L}$ that gets quoted in this context is not a rule of thumb and not
an empirical finding. It is where the derivative of $\frac{L}{k} + k$
vanishes, which is Program~\ref{prog:F11}'s stationary point used on a cost
rather than on a loss. The script asserts the numerical minimum agrees with
$\sqrt{L}$ rather than asserting a tolerance.
\end{note}
\end{fr}

\begin{fr}
\begin{aibox}
That factor is the difference between a model that fits on a card and one that
does not, and the price is stated exactly: one extra forward pass, which is a
third of a training step rather than a doubling of it, because a step is a
forward and a backward and the backward is the expensive half.

It is also why the knob is usually on or off rather than tuned. The optimum is
$\sqrt{L}$ and the curve near a minimum is flat \dash{} which
Program~\ref{prog:P17} will say properly \dash{} so anything within a factor
of two of $\sqrt{L}$ costs almost the same.
\end{aibox}
\end{fr}

\begin{fr}
One more thing this section owes, and it is the tool everyone reaches for when
a gradient looks wrong.

You can always estimate a derivative with a finite difference, as
Program~\ref{prog:F11} did. Why is that a test of an implementation and not an
implementation itself?
\dotline
\nextframe
\end{fr}

\begin{fr}
\begin{ansblock}
Because it costs one evaluation per input, which is forward mode's cost with
none of its accuracy \dash{} and Program~\ref{prog:F11} measured that its
error has a floor no step size can get under.
\end{ansblock}

Two separate objections, and the second is the one people forget. F11's
U-curve showed the error falling as the step shrinks and then \emph{rising}
again as cancellation takes over, with a best case that is nowhere near
machine precision.

\textbf{So a gradient check that agrees to four or five digits is what a
correct implementation looks like}, and demanding fifteen would be demanding
that the test be better than it can be.

This program's own differentiator is checked exactly that way: against a
central difference at $\val{p16.check.n}$ inputs, agreeing to better than
$\val{p16.check.bound}$. That is the check passing, not the check being weak.
\end{fr}

% ============================================================
\section{Three times it answers a different question}
\label{sec:P16-silent}
\index{autodiff!silent failures}

\begin{fr}
Everything so far has been about what autodiff computes. This section is about
three situations in which it computes something else and says nothing, which
is the failure mode that costs the most time because there is no error to
search for.

All three are demonstrated on a reverse-mode differentiator of about forty
lines, written for this program, so nothing here depends on a claim about any
particular framework.

The first. $\relu$ is not differentiable at zero. What should
$\frac{\mathrm{d}}{\mathrm{d}x}\relu(x)$ return at exactly $x = 0$?
\dotline
\nextframe
\end{fr}

\begin{fr}
\begin{ansblock}
There is no right answer. The limit from below is
$\val{p16.relu.left}$ and from above is $\val{p16.relu.right}$, so the
derivative does not exist and an implementation must \emph{choose}.
\end{ansblock}

And it chooses without telling you. This one returns
$\val{p16.relu.chosen}$; a different one may return the other; and nothing in
either case raises anything.

\transcript{p16-silent}

\begin{trapbox}
\textbf{\enquote{The gradient is wrong at zero} is the wrong complaint.} There
is no gradient at zero to be wrong about. What is true is that the number you
get is a convention, and that two implementations may disagree about it while
both being defensible.

In practice it almost never matters, because hitting exactly zero in floating
point is rare \dash{} and \enquote{almost never} is doing real work in that
sentence. A layer whose inputs are all exactly zero, or a mask, or a padded
batch, will hit it every single time.
\end{trapbox}
\end{fr}

\begin{fr}
The second, and it is the one that produces a model that trains and does not
learn.

Somewhere in the computation a value is detached from the graph \dash{} copied
out and put back in without its history. Everything downstream still computes
correctly. What happens to the gradient?
\nextframe
\end{fr}

\begin{fr}
\begin{ansblock}
The path through that value contributes nothing. The gradient is not wrong in
the sense of miscalculated; a term is simply \emph{missing}.
\end{ansblock}

Measured on the same forty-line differentiator: detaching one intermediate
leaves the loss value \emph{identical to the last bit} and changes one
component of the gradient by $\val{p16.detach.gap}$.

\textbf{The loss printed in the log is exactly the same number.} Every
diagnostic that looks at the forward pass agrees with a correct run. The only
evidence is that a part of the model stops improving, and the usual first
guess is the learning rate.

\begin{note}
Detaching is often deliberate and correct \dash{} a target that should not be
trained through, a value used for logging, a stop-gradient in an algorithm
that requires one. The hazard is not the operation, it is that its effect is
invisible in every place people look.
\end{note}
\end{fr}

\begin{fr}
The third. A value that the tape still refers to is overwritten in place after
being used.

Say what the backward pass computes with, before reading on.
\dotline
\nextframe
\end{fr}

\begin{fr}
\begin{ansblock}
Whichever value the tape recorded at the time \dash{} which may be the old one
or the new one, depending on whether the local derivative was stored or the
tensor was.
\end{ansblock}

In the listing above, the product of $\val{p16.inplace.a}$ and
$\val{p16.inplace.used}$ stored the second of them as the first one's local
derivative. Overwriting that second value afterwards leaves the gradient
computed with $\val{p16.inplace.used}$ while the value now holds
$\val{p16.inplace.now}$.

\textbf{So the gradient is consistent with a computation that no longer
exists.} It is not garbage, which is what makes it dangerous \dash{} it is a
perfectly reasonable number for a different problem.

\begin{note}
This is the one of the three that real frameworks do usually catch, because
they version their buffers and can tell that one changed after being recorded.
Whether a particular one catches a particular case is a fact about a version,
which is why no framework is named anywhere in this program.

What is durable is the shape of the hazard: \textbf{a tape stores what it saw,
and a value that changes afterwards makes the tape a record of something that
did not happen.}
\end{note}
\end{fr}

\begin{fr}
The three have one thing in common, and naming it is more useful than the
three individually. What is it?
\dotline
\nextframe
\end{fr}

\begin{fr}
\begin{ansblock}
In all three the forward pass is correct and the loss is exactly right. Only
the derivative is not what you meant, and nothing looks at the derivative.
\end{ansblock}

Which is why the debugging habit that works is not reading the training curve.
It is Program~\ref{prog:F11}'s finite difference, used as the test this
program said it is: pick a few parameters, perturb them, and check the number
the framework gave you against the number the definition gives.

\textbf{That is the only instrument that looks at the derivative directly},
and it takes a few lines. All three of this section's failures show up in it
immediately, and none of them shows up anywhere else.
\end{fr}

\begin{fr}
One closing sentence, because the program has been about a mechanism and it is
worth saying what the mechanism is not.

Automatic differentiation is not numerical and it is not symbolic. It does not
approximate anything, so it has none of the finite difference's error floor;
and it never builds an expression for the derivative, so it has none of the
blow-up that differentiating a deep expression symbolically would produce.

\textbf{It evaluates the chain rule on the graph the program actually ran},
one edge at a time, to the same precision as the forward pass. That is all
$\code{backward()}$ is, and there is nothing else in it.
\end{fr}

% ============================================================
\begin{summarybox}
\sumitem{1--2}{\result{\textbf{The derivative of a function from $m$ inputs to $n$
  outputs is an $n \times m$ matrix} \dash{} the Jacobian, one row per output
  and one column per input}.}
\sumitem{3--4}{\result{\textbf{Program~\ref{prog:P15}'s gradient is the one-output
  case}, a single row, so nothing new was defined here}.}
\sumitem{5--6}{\result{\textbf{A linear map is its own Jacobian}, which is
  Program~\ref{prog:F11}'s statement about a straight line one dimension count
  up. An elementwise activation's Jacobian is diagonal}.}
\sumitem{9--10}{\result{\textbf{The chain rule multiplies Jacobians}, in the order the
  shapes allow, which is Program~\ref{prog:P06}'s composition}.}
\sumitem{11--14}{\result{\textbf{Nobody forms a Jacobian.} What is computed is a
  product with one, and the intermediates are then vectors rather than
  matrices}.}
\sumitem{14--15}{\result{The same product bracketed from either end costs differently
  \dash{} Program~\ref{prog:P06} measured $\val{p06.cost.left}$ against
  $\val{p06.cost.right}$ for an identical answer. \textbf{Those two brackets
  are the two modes.}}}
\sumitem{16--17}{\result{\textbf{Forward mode costs one pass per input; reverse costs
  one per output.} That is the whole decision procedure and it is two
  integers}.}
\sumitem{17--18}{\result{\textbf{A scalar loss makes reverse mode cheaper by a factor
  of the parameter count.} Nothing else about deep learning enters the
  argument}.}
\sumitem{19--20}{\result{Measured exactly over the rationals: $\val{p16.fwd.muls}$
  multiplications forward against $\val{p16.rev.muls}$ reverse for an
  identical gradient, and \textbf{the ratio is the input count}, swept over
  four widths}.}
\sumitem{21--22}{\result{Program~\ref{prog:P13} is what makes the backward walk well
  defined: the graph has a topological order and \textbf{every one of them
  computes the same values}}.}
\sumitem{23--24}{\result{\textbf{Reverse mode must keep every layer's input}, because
  the Jacobian's entries depend on it. Memory grows with depth and with batch
  size; forward mode needs none of it}.}
\sumitem{25--26}{\result{\textbf{Checkpointing every $k$th layer holds
  $\frac{L}{k} + k$ activations}, minimised at $k = \sqrt{L}$ \dash{} which is
  where a derivative vanishes, not a rule of thumb}.}
\sumitem{26--27}{\result{At $L = \val{p16.ckpt.layers}$ that is
  $\val{p16.ckpt.peak}$ held instead of $\val{p16.ckpt.layers}$, a factor of
  $\val{p16.ckpt.saving}$, for exactly $\val{p16.ckpt.passes}$ forward passes
  rather than one}.}
\sumitem{28--29}{\result{\textbf{A finite difference is a test, not an
  implementation}: it costs one evaluation per input and
  Program~\ref{prog:F11} measured its error floor, so agreement to a handful
  of digits is what correct looks like}.}
\sumitem{30--31}{\result{\textbf{$\relu$ at exactly zero has no derivative}, so an
  implementation chooses \dash{} the limits are $\val{p16.relu.left}$ and
  $\val{p16.relu.right}$ and this one returns $\val{p16.relu.chosen}$, without
  saying so}.}
\sumitem{32--33}{\result{\textbf{A detached value leaves the loss identical to the
  last bit} and a gradient component different by $\val{p16.detach.gap}$. A
  term is missing rather than miscalculated}.}
\sumitem{34--35}{\result{\textbf{An in-place write makes the tape a record of a
  computation that did not happen}: the gradient here uses
  $\val{p16.inplace.used}$ while the value holds $\val{p16.inplace.now}$}.}
\sumitem{36--37}{\result{\textbf{All three leave the forward pass exactly right}, so
  only a direct check of the derivative finds them \dash{} which is what the
  finite difference is for}.}
\sumitem{38}{\result{Autodiff is neither numerical nor symbolic. \textbf{It evaluates
  the chain rule on the graph the program actually ran}, to the same precision
  as the forward pass}.}
\end{summarybox}

\canyou

\begin{testexercises}
\item A function takes $4$ inputs and returns $7$ outputs. Give the shape of
  its Jacobian, and say how many forward-mode passes and how many reverse-mode
  passes it would take to fill the whole thing in.
  \answerto{$7 \times 4$. Forward mode needs $4$ passes, one per input;
  reverse needs $7$, one per output. Here forward is cheaper, which is the
  opposite of training and for exactly the same reason \dash{} the shape.}
\item Why is the Jacobian of an elementwise activation diagonal, and what does
  that buy?
  \answerto{Because output $i$ depends on input $i$ alone, so every
  off-diagonal partial derivative is zero. It buys everything: multiplying by
  a diagonal matrix is one multiplication per entry rather than $n$, so an
  activation's contribution to a backward pass costs the same as the
  activation itself.}
\item A stack is $L = 400$ layers deep. Give the best checkpoint spacing, the
  peak number of activations held, and the compute cost.
  \answerto{$k = \sqrt{400} = 20$, so $20$ checkpoints plus a segment of $20$
  is $40$ held against $400$ \dash{} a factor of ten \dash{} for exactly two
  forward passes rather than one. The factor is $\sqrt{L}/2$ in general, so it
  improves as the stack gets deeper, which is why the technique is used where
  it is.}
\item Your gradient check agrees to six digits for every parameter except one,
  where it disagrees in the first digit. What are the three candidates from
  section 5, and how would you tell them apart?
  \answerto{A non-differentiable point, a detached path, or an in-place write.
  Perturb that one parameter and watch: if the disagreement vanishes when you
  move slightly off the point, it is the first; if the analytic gradient is
  exactly zero where the numerical one is not, it is a missing path; if the
  forward value changes when you reorder unrelated operations, it is the
  third.}
\item Why does the argument for reverse mode not mention the network's depth?
  \answerto{Because depth affects both modes identically \dash{} each is one
  walk along the chain, whichever end it starts from. What differs is how many
  walks are needed, and that is set by the number of inputs against the number
  of outputs. Depth changes the cost of a pass; the shape changes the number
  of passes.}
\end{testexercises}

\begin{furtherproblems}
\item Show that if $f$ is linear then forward and reverse mode compute the
  same thing in the same number of operations when there is one input and one
  output.
  \answerto{With one input and one output the Jacobian is a single number, the
  chain is a product of scalars, and both bracketings are the same product in
  a different order \dash{} the same count either way. The two modes only
  separate when a dimension is bigger than one, which is another way of saying
  the shape is the whole argument.}
\item A model has $P$ parameters and the loss is a vector of $C$ per-class
  losses rather than a scalar. What does that do to the cost of getting every
  derivative?
  \answerto{It multiplies reverse mode's cost by $C$, because reverse needs
  one pass per output. That is why a training loop reduces to a scalar before
  differentiating: summing or averaging the per-class losses first turns $C$
  passes into one, and it is the same reduction that makes the whole argument
  in section 3 work.}
\item Forward mode needs no stored activations. Why is that not an argument
  for using it in training?
  \answerto{Because it would need one pass per parameter, and the parameter
  count is the largest number in the problem. Trading a memory cost that grows
  with depth for a compute cost that grows with the model size is not a trade
  anyone would take. It does make forward mode the right tool for a
  derivative with respect to one or two things, which is where it is used.}
\item Section 4 says checkpointing costs exactly two forward passes. Under
  what recomputation scheme would it cost more, and why is that scheme
  sometimes chosen anyway?
  \answerto{Recomputing recursively \dash{} checkpointing inside the segments
  as well \dash{} costs more passes and less memory still, roughly
  $O(\mfalogplain L)$ passes for $O(\mfalogplain L)$ storage. It is chosen when memory is the
  binding constraint and compute is not, which is the same trade one level
  further along, and the arithmetic is the same minimisation done twice.}
\item Why can a gradient be exactly zero for a legitimate reason and for a bug
  that looks identical?
  \answerto{Legitimately, at a stationary point, past a $\relu$ that is off, or
  where a parameter genuinely does not affect the loss. As a bug, from a
  detached path, which contributes nothing by construction. The two are
  indistinguishable from the gradient alone, and the way to tell is to
  perturb the parameter and see whether the loss moves: if it moves and the
  gradient is zero, the path is broken rather than flat.}
\end{furtherproblems}
